\documentclass[12pt]{article}
\usepackage{statproj}

\definecolor{myblue}{HTML}{4c72b0}
\definecolor{myorange}{HTML}{dd8452}
\definecolor{mygreen}{HTML}{55a868}

\title{A Log-Linear Regression Analysis of Allometric Scaling and
Species Effects on Penguin Body Mass}
\author{Avigyan Chakraborty}
\date{April, 2026}

\begin{document}

\maketitlepage{\thetitle}{\theauthor}{\thedate}{%
  placeholder for abstract%
}

% TODO: change the red color in this
\makecontents

\section{Introduction}

insert introduction here.

\section{Data Description}

The dataset used in this study is the LIDC-IDRI (Lung Image
Database Consortium and Image Database Resource Initiative)
database\cite{armato2011lidc}.
It contains thoracic CT scans in which up to four radiologists
outlined each lung nodule of at least $3$\,mm and rated its
appearance and its likely malignancy.

The unit of analysis is the \emph{nodule}. Geometric features are
averaged across a nodule's readers and semantic ratings are
summarised by their median. We use the median because the ratings
are ordinal, so the gaps between grades have no fixed numerical
meaning, and because it is less sensitive to a single outlying
reader than the mean.

The variables used in this study are as follows.

\paragraph{Response}
\begin{itemize}
    \ii \textbf{Malignancy (\texttt{mal\_median})}:
    The median of the readers' ratings, on an ordered scale from
    $1$ (highly unlikely to be malignant) to $5$
    (highly suspicious). A median falling between two grades is
    rounded up, which affects $548$ nodules ($20.8\%$).
    Each reader's rating is also retained for the latent-grade model.
    Discrete (ordinal), with support $\{1, 2, 3, 4, 5\}$.
\end{itemize}

\paragraph{Primary model: geometric features}
These are computed from the outlines drawn by the radiologists and
are treated as continuous random variables.
\begin{itemize}
    \ii \textbf{Log Volume (\texttt{log\_volume})}:
    The natural logarithm of the nodule's volume in mm$^3$.
    The logarithm is used because volume is strongly right-skewed.
    Continuous, with support $\mathbb{R}$.
    \ii \textbf{Elongation (\texttt{elongation})}:
    A value near $1$ means that the nodule is not elongated.
    Continuous, with support $[0, 1]$.
    \ii \textbf{Flatness (\texttt{flatness})}:
    A value near $1$ means that the nodule is not flat.
    Continuous, with support $[0, 1]$ and a point mass at $0$.
    \ii \textbf{Solidity (\texttt{solidity})}:
    The number of voxels in the nodule divided by the number in its
    convex hull. A low value indicates an irregular outline.
    Continuous, with support $(0, 1]$.
    \ii \textbf{Compactness (\texttt{compactness})}:
    $V^{2/3}/S$, for volume $V$ and surface area $S$.
    A high value means little surface area relative to volume.
    Continuous, with support $(0, \infty)$.
\end{itemize}

\paragraph{Secondary model: primary features and semantic ratings}
These are ordinal ratings from $1$ to $5$, given by the same
radiologists who rate malignancy.
The secondary model adds them to the five
geometric features above. Since each is a median across readers, all
five are discrete (ordinal), with support
$\{1, 1.5, 2, \dots, 4.5, 5\}$.
\begin{itemize}
    \ii \textbf{Spiculation (\texttt{spiculation})}:
    Spiky extensions from the edge, from none to marked.
    \ii \textbf{Lobulation (\texttt{lobulation})}:
    Rounded bulges along the edge, from none to marked.
    \ii \textbf{Margin (\texttt{margin})}:
    Definition of the edge, from poorly defined to sharp.
    \ii \textbf{Texture (\texttt{texture})}:
    Internal density, from non-solid to solid.
    \ii \textbf{Sphericity (\texttt{sphericity})}:
    Overall roundness, from linear to round.
\end{itemize}

The data are split by patient into a training set ($609$ patients,
$1809$ nodules) and a test set ($261$ patients, $828$ nodules).
Flatness is exactly $0$ for $45$ training nodules whose outlines lie
in a single plane.

\subsection{Summary Statistics}

All summary statistics are computed on the training set.
\Cref{table:summary_all} gives the unstandardised values of the
predictors, and \cref{table:summary_grades} breaks them down by
median malignancy grade.

\begin{table}[H]
  \centering
  \begin{tabular}{|| c | c | c | c | c | c ||}
    \hline
    \textbf{Model}
    & \textbf{Variable} & \textbf{Mean}
    & \textbf{Minimum}
    & \textbf{Maximum} & \textbf{Std. Dev.} \\[0.5ex]
    \hline\hline
    \multirow{5}{*}{\textbf{Primary}}
    & \textbf{Log Volume} & 5.118 & 2.442 & 10.345 & 1.394 \\
    & \textbf{Elongation} & 0.753 & 0.306 & 0.988 & 0.121 \\
    & \textbf{Flatness} & 0.555 & 0.000 & 0.921 & 0.164 \\
    & \textbf{Solidity} & 0.795 & 0.228 & 1.000 & 0.106 \\
    & \textbf{Compactness} & 0.177 & 0.058 & 0.312 & 0.044 \\
    \hline
    \multirow{5}{*}{\textbf{Secondary}}
    & \textbf{Spiculation} & 1.476 & 1.000 & 5.000 & 0.889 \\
    & \textbf{Lobulation} & 1.572 & 1.000 & 5.000 & 0.883 \\
    & \textbf{Margin} & 3.899 & 1.000 & 5.000 & 1.138 \\
    & \textbf{Texture} & 4.363 & 1.000 & 5.000 & 1.202 \\
    & \textbf{Sphericity} & 3.789 & 1.000 & 5.000 & 0.826 \\
    \hline
  \end{tabular}
  \caption{Overall Summary Statistics ($n = 1809$ training nodules)}
  \label{table:summary_all}
\end{table}

\begin{table}[H]
  \centering
  \resizebox{\textwidth}{!}{%
    \begin{tabular}{|| c | c | c | c | c | c | c ||}
      \hline
      \textbf{Model} & \textbf{Variable}
      &
      \begin{tabular}{@{}c@{}}\textbf{Grade 1} \\ \textbf{(233)}
      \end{tabular}
      &
      \begin{tabular}{@{}c@{}}\textbf{Grade 2} \\ \textbf{(378)}
      \end{tabular}
      &
      \begin{tabular}{@{}c@{}}\textbf{Grade 3} \\ \textbf{(823)}
      \end{tabular}
      &
      \begin{tabular}{@{}c@{}}\textbf{Grade 4} \\ \textbf{(239)}
      \end{tabular}
      &
      \begin{tabular}{@{}c@{}}\textbf{Grade 5} \\ \textbf{(136)}
      \end{tabular} \\[0.5ex]
      \hline\hline
      \multirow{5}{*}{\textbf{Primary}}
      & \textbf{Log Volume} & 4.563 (0.981) & 4.474 (0.935) & 4.779
      (0.983) & 6.349 (1.269) & 7.744 (1.082) \\
      & \textbf{Elongation} & 0.749 (0.135) & 0.731 (0.125) & 0.758
      (0.117) & 0.773 (0.115) & 0.761 (0.109) \\
      & \textbf{Flatness} & 0.591 (0.150) & 0.511 (0.173) & 0.548
      (0.167) & 0.584 (0.151) & 0.602 (0.132) \\
      & \textbf{Solidity} & 0.842 (0.092) & 0.812 (0.092) & 0.808
      (0.093) & 0.742 (0.104) & 0.676 (0.125) \\
      & \textbf{Compactness} & 0.184 (0.049) & 0.187 (0.043) & 0.183
      (0.040) & 0.154 (0.037) & 0.135 (0.031) \\
      \hline
      \multirow{5}{*}{\textbf{Secondary}}
      & \textbf{Spiculation} & 1.172 (0.645) & 1.200 (0.632) & 1.394
      (0.766) & 1.795 (0.930) & 2.699 (1.243) \\
      & \textbf{Lobulation} & 1.255 (0.700) & 1.300 (0.728) & 1.509
      (0.819) & 1.933 (0.878) & 2.618 (0.970) \\
      & \textbf{Margin} & 4.674 (0.793) & 3.971 (1.148) & 3.756
      (1.159) & 3.669 (1.076) & 3.643 (1.011) \\
      & \textbf{Texture} & 4.880 (0.560) & 4.344 (1.253) & 4.197
      (1.323) & 4.324 (1.139) & 4.603 (0.893) \\
      & \textbf{Sphericity} & 4.043 (0.924) & 3.810 (0.877) & 3.779
      (0.781) & 3.653 (0.790) & 3.599 (0.735) \\
      \hline
  \end{tabular}}
  \caption{Mean (Std. Dev.) by Median Malignancy Grade
  (number of training nodules in brackets)}
  \label{table:summary_grades}
\end{table}

\subsection{Data Source}

The dataset used in this study is the LIDC-IDRI
database\cite{armato2011lidc}, hosted by The Cancer Imaging
Archive\cite{clark2013tcia} at
\href{https://www.cancerimagingarchive.net/collection/lidc-idri/}
{LIDC-IDRI (TCIA)}.
The full image collection is about $133$\,GB, but the radiologists'
annotations are bundled, already parsed, with the Python library
\texttt{pylidc}\cite{hancock2016pylidc}. Every quantity used in
this project is computed from these annotations with
\texttt{pylidc} version $0.2.3$, so no CT images were downloaded.

\section{Model Specification}

\subsection{Logistic Regression Model (Naive Proportional Odds)}

\paragraph{Methodology}
Let $Y_i \in \{1, \dots, 5\}$ be the median malignancy grade of nodule
$i$, and let $x_i$ be its vector of predictors. Since $Y_i$ is ordered,
we model its cumulative probabilities with a proportional odds model
(an ordinal logistic regression):
\[
  \operatorname{logit} P(Y_i \le k \mid x_i)
  = \alpha_k - x_i\trans\beta, \qquad k = 1, \dots, 4,
\]
where $\alpha_1 < \dots < \alpha_4$ are thresholds. A positive
$\beta_j$ means that a larger $x_{ij}$ shifts a nodule towards higher
grades. The odds of being above any given grade are multiplied by
$e^{\beta_j}$ for each one-unit increase in $x_{ij}$, and this factor
is assumed to be the same at every threshold. The model is
\emph{naive} because it takes the median as the true grade and ignores
disagreement between readers.

Writing $F(t) = e^t / (1 + e^t)$ for the logistic distribution
function, with $\alpha_0 = -\infty$ and $\alpha_5 = \infty$, and
treating nodules as independent, the likelihood is
\[
  L(\alpha, \beta)
  = \prod_{i=1}^{n} \prod_{k=1}^{5}
  \bigl[ F(\alpha_k - x_i\trans\beta)
  - F(\alpha_{k-1} - x_i\trans\beta) \bigr]^{\mathbf{1}\{Y_i = k\}}.
\]

The models are fitted by maximising $L$ on the training set with
\texttt{polr}. All predictors are standardised, so one unit is
one standard deviation. One nodule with missing solidity is excluded,
leaving $n = 1808$ for both models. Confidence intervals for the odds
ratios are profile-likelihood intervals, obtained by inverting the
likelihood ratio test of $H_0 \colon \beta_j = b$ using the
$\chi^2_1$ limit given by Wilks' theorem, and then exponentiating
the endpoints.

\paragraph{Primary Model}
The fitted model is
\begin{align*}
  \operatorname{logit} \hat{P}(Y_i \le k \mid x_i) = \hat{\alpha}_k
  - \bigl(&1.330\cdot\texttt{log\_volume}
    + 0.333\cdot\texttt{elongation} \\
    &- 0.173\cdot\texttt{flatness}
    - 0.342\cdot\texttt{solidity} \\
  &+ 0.047\cdot\texttt{compactness}\bigr),
\end{align*}
with $\hat{\alpha} = (-2.503, -1.108, 1.674, 3.488)$ and maximised
log-likelihood $-2155.134$. The estimates are given in
\cref{table:polr_primary}.

\begin{table}[H]
  \centering
  \begin{tabular}{|| c | c | c | c | c ||}
    \hline
    \textbf{Predictor}
    & $\hat{\beta}$
    & \textbf{Std. Error}
    & \textbf{Odds Ratio}
    & \textbf{$95\%$ CI} \\[0.5ex]
    \hline\hline
    \texttt{log\_volume} & 1.330 & 0.065 & 3.781 & (3.335, 4.296) \\
    \texttt{elongation} & 0.333 & 0.055 & 1.395 & (1.253, 1.555) \\
    \texttt{flatness} & -0.173 & 0.053 & 0.841 & (0.757, 0.934) \\
    \texttt{solidity} & -0.342 & 0.065 & 0.710 & (0.625, 0.806) \\
    \texttt{compactness} & 0.047 & 0.059 & 1.048 & (0.934, 1.175) \\
    \hline
  \end{tabular}
  \caption{Primary Proportional Odds Model}
  \label{table:polr_primary}
\end{table}

\paragraph{Interpretation}
Each odds ratio is for a one standard deviation increase in that
predictor, holding the others constant.
\begin{itemize}
    \ii \textbf{Size:} Log volume has by far the largest effect. One
    standard deviation, roughly a fourfold increase in volume,
    multiplies the odds of a higher malignancy grade by $3.78$.
    \ii \textbf{Solidity:} A more solid nodule has lower odds of a
    higher grade (odds ratio $0.71$), so irregular outlines are rated
    as more suspicious.
    \ii \textbf{Elongation and Flatness:} Elongation is close to $1$
    for a nodule that is not elongated, so an odds ratio of $1.40$
    means that more elongated nodules have lower odds of a higher
    grade. Flatness works in the other direction: its odds ratio of
    $0.84$ means that flatter nodules have higher odds.
    \ii \textbf{Compactness:} The confidence interval contains $1$, so
    the data are compatible with no effect once the other predictors
    are included.
\end{itemize}

\paragraph{Secondary Model}
The geometric features only describe the outline of a nodule, whereas
radiologists also judge its appearance on the scan. We therefore add
the semantic ratings to measure how much this assessment explains
beyond geometry: spiculation, lobulation and margin describe the edge
in terms that radiologists themselves use, texture records internal
density, which the outlines cannot capture, and sphericity gives a
reader's view of overall shape. Adding these ratings gives
\begin{align*}
  \operatorname{logit} \hat{P}(Y_i \le k \mid x_i) = \hat{\alpha}_k
  - \bigl(&1.210\cdot\texttt{log\_volume}
    + 0.321\cdot\texttt{elongation} \\
    &- 0.065\cdot\texttt{flatness}
    - 0.160\cdot\texttt{solidity} \\
    &+ 0.002\cdot\texttt{compactness}
    + 0.346\cdot\texttt{spiculation} \\
    &+ 0.091\cdot\texttt{lobulation}
    - 0.280\cdot\texttt{margin} \\
    &+ 0.063\cdot\texttt{texture}
  + 0.002\cdot\texttt{sphericity}\bigr),
\end{align*}
with $\hat{\alpha} = (-2.572, -1.144, 1.735, 3.617)$ and maximised
log-likelihood $-2108.642$. The estimates are given in
\cref{table:polr_secondary}.

\begin{table}[H]
  \centering
  \begin{tabular}{|| c | c | c | c | c ||}
    \hline
    \textbf{Predictor}
    & $\hat{\beta}$
    & \textbf{Std. Error}
    & \textbf{Odds Ratio}
    & \textbf{$95\%$ CI} \\[0.5ex]
    \hline\hline
    \texttt{log\_volume} & 1.210 & 0.067 & 3.352 & (2.942, 3.828) \\
    \texttt{elongation} & 0.321 & 0.056 & 1.379 & (1.237, 1.539) \\
    \texttt{flatness} & -0.065 & 0.056 & 0.937 & (0.840, 1.045) \\
    \texttt{solidity} & -0.160 & 0.069 & 0.852 & (0.744, 0.976) \\
    \texttt{compactness} & 0.002 & 0.059 & 1.002 & (0.893, 1.126) \\
    \texttt{spiculation} & 0.346 & 0.067 & 1.413 & (1.241, 1.613) \\
    \texttt{lobulation} & 0.091 & 0.066 & 1.096 & (0.964, 1.246) \\
    \texttt{margin} & -0.280 & 0.074 & 0.756 & (0.654, 0.874) \\
    \texttt{texture} & 0.063 & 0.067 & 1.065 & (0.933, 1.215) \\
    \texttt{sphericity} & 0.002 & 0.051 & 1.002 & (0.907, 1.107) \\
    \hline
  \end{tabular}
  \caption{Secondary Proportional Odds Model}
  \label{table:polr_secondary}
\end{table}

\paragraph{Interpretation}
\begin{itemize}
    \ii \textbf{Semantic Ratings:} More spiculated nodules have higher
    odds of a higher grade (odds ratio $1.41$), and nodules with
    sharper margins have lower odds ($0.76$). The confidence intervals
    for lobulation, texture and sphericity contain $1$.
    \ii \textbf{Geometric Features:} Size remains the largest effect
    (odds ratio $3.35$), and elongation barely changes ($1.38$). The
    solidity effect moves towards $1$ ($0.85$), and the interval for
    flatness now contains $1$, which suggests that the ratings capture
    some of the same information about shape.
    \ii \textbf{Circularity:} The semantic ratings come from the same
    radiologists who rate malignancy, so these associations may partly
    reflect a reader's overall impression of a nodule. The primary
    model therefore remains our main result.
\end{itemize}

\section{Statistical Inference}

\subsection{Confidence Intervals}

In multiple linear regression,
estimating the point values for the coefficients
$(\hat{\beta}_i)$ provides a single best guess for the
population parameters. However,
to account for sampling variability and quantify the
precision of these estimates, we construct $95\%$
confidence intervals for each predictor, the mean response,
and future individual observations.

\subsubsection{Coefficient Intervals}
The true parameter $(\hat{\beta}_i)$ is estimated to fall
within the following interval with a $95\%$ level of confidence:
\[
  \hat{\beta}_i \pm t_{{\alpha}/{2}}(n-p-1)\cdot SE(\hat{\beta}_i)
.\]
Plugging in the values, we derive the following interval estimates:
\begin{table}[H]
  \centering
  \begin{tabular}{|| c | c | c c ||}
    \hline
    \textbf{Predictor Variable}
    & \textbf{Estimate}
    & \textbf{Lower Bound ($2.5\%$)}
    & \textbf{Upper Bound ($97.5\%$)} \\[0.5ex]
    \hline\hline
    \texttt{Intercept} ($\hat{\beta}_0$) & 0.0759 & -1.2559 & 1.4077 \\
    $\log(\texttt{Flipper\_Length})$ ($\hat{\beta}_1$) & 0.9036 &
    0.6099 & 1.1972 \\
    $\log(\texttt{Culmen\_Length})$ ($\hat{\beta}_2$) & 0.4364 &
    0.2864 & 0.5865 \\
    $\log(\texttt{Culmen\_Depth})$ ($\hat{\beta}_3$) & 0.6181 &
    0.4598 & 0.7763 \\
    \texttt{Is\_Gen} ($\hat{\beta}_4$) & 0.2341 & 0.1648 & 0.3034 \\
    \texttt{Is\_Chin} ($\hat{\beta}_5$) & -0.1195 & -0.1592 & -0.0798 \\
    \hline
  \end{tabular}
  \caption{$95\%$ Confidence Intervals for the Regression Parameters}
\end{table}

Because both the response variable and the continuous predictors were
log-transformed, the coefficients are interpreted as
elasticities, representing the expected percentage change in
body mass for a $1\%$ increase in the predictor,
holding all other variables constant.
\begin{itemize}
    \ii \textbf{Flipper Length:} We are $95\%$ confident that for every $1\%$
    increase in a penguin's flipper length,
    its expected body mass increases by between $0.61\%$ and $1.20\%$.

    \ii \textbf{Culmen Length:} We are $95\%$ confident that a $1\%$
    increase in culmen length corresponds to an expected body mass
    increase of between $0.29\%$ and $0.59\%$.

    \ii \textbf{Culmen Depth:} We are $95\%$ confident that a $1\%$
    increase in culmen depth is associated with a
    $0.46\%$ to $0.78\%$ increase in expected body mass.

    \ii \textbf{Species Effects (\texttt{Is\_Gen}):}
    We are $95\%$ confident that,
    holding all physical measurements constant,
    a Gentoo penguin will weigh between $17.9\%$ and $35.4\%$
    more than the baseline Adelie species.

    \ii \textbf{Species Effects (\texttt{Is\_Chin}):}
    We are $95\%$ confident that,
    holding all physical measurements constant,
    a Chinstrap penguin will weigh between $7.7\%$ and $14.7\%$
    less than the Adelie species.

    \ii \textbf{Intercept:} The $95\%$ confidence interval
    for the intercept ranges from $-1.2559$ to $1.4077$.
    The intercept corresponds to the expected
    log body mass of a baseline Adelie penguin
    when all log-transformed predictors are equal to zero.
    This implies that the original
    measurements are equal to 1 (since $\log(1) = 0$),
    which is not biologically meaningful.
    Thus, the intercept has limited practical
    interpretation in this context.
\end{itemize}

\subsubsection{Confidence Intervals for the Mean Response}
Let $x_0$ be
a column vector containing a specific set of predictor
values. To estimate the expected value $\mathbb{E}[\log(\texttt{Body\_M})]$ at
$x_0$, we utilize the matrix $X$
to account for the covariance among predictors. The $95\%$
confidence interval for the mean response is given by:
\[
  \hat{y}_0 \pm t_{\alpha/2}(n-p-1) \cdot
  \sqrt{\text{MSE} \cdot x_0^T
  (X^T X)^{-1} x_0}
.\]
Here, MSE is the Mean Squared Error, estimating the error
variance $\sigma^2$.

\paragraph{Application: Expected Mass Across Species}
To isolate the structural effects of the categorical
variables, we construct a theoretical standardized penguin
utilizing the mean of the log-transformed morphological
measurements. We then predict the expected $\log(\texttt{Body\_M})$
across all three species profiles, keeping these baseline
log-dimensions mathematically constant.

\begin{table}[H]
  \centering
  \begin{tabular}{|| c | c | c c ||}
    \hline
    \textbf{Species}
    & \textbf{Point Estimate}
    & \textbf{Lower Bound}
    & \textbf{Upper Bound} \\
    \hline\hline
    \textbf{Adelie (Baseline)} & 8.2651 & 8.2337 & 8.2964 \\
    \textbf{Chinstrap} & 8.1456 & 8.1178 & 8.1733 \\
    \textbf{Gentoo} & 8.4992 & 8.4581 & 8.5402 \\
    \hline
  \end{tabular}
  \caption{$95\%$ Confidence Intervals for Mean Response}
\end{table}

Observing the Chinstrap and Gentoo confidence bounds,
we see parallel shifts relative to the Adelie baseline.
Because the input traits are identical, we are $95\%$
confident that the average log body mass of a Gentoo
population will always remain strictly higher than an
identically proportioned Adelie population, whereas a
Chinstrap population will remain strictly lower.

\subsubsection{Prediction Intervals}
While the confidence interval predicts
$\mathbb{E}[\log(\texttt{Body\_M})]$
for all penguins with a specific set of traits, the
prediction interval predicts the exact
$\log(\texttt{Body\_M})$
of a
single, individual new penguin. Because individual
observations naturally vary around the population mean,
predicting a single point carries additional irreducible
uncertainty. The $95\%$ prediction interval
of $\log(\texttt{Body\_M})$ is mathematically
wider and is expressed as:
\[
  \hat{y}_0 \pm t_{\alpha/2}(n-p-1) \cdot
  \sqrt{\text{MSE} \cdot \left(1 + x_0^T
  (X^T X)^{-1} x_0\right)}
.\]
The addition of $1$ inside the parenthetical term accounts
for the variance of the residual error ($\sigma^2$)
associated with drawing a new, random individual from the
population.

\paragraph{Application: Predicting Individual Penguins}
Using the same theoretical baseline log-dimensions, we now
calculate the intervals for randomly sampling a single new
penguin from each respective species.

\begin{table}[H]
  \centering
  \begin{tabular}{|| c | c | c c ||}
    \hline
    \textbf{Species}
    & \textbf{Point Estimate}
    & \textbf{Lower Bound}
    & \textbf{Upper Bound} \\
    \hline\hline
    \textbf{Adelie (Baseline)} & 8.2651 & 8.1107 & 8.4195 \\
    \textbf{Chinstrap} & 8.1456 & 7.9918 & 8.2993 \\
    \textbf{Gentoo} & 8.4992 & 8.3425 & 8.6558 \\
    \hline
  \end{tabular}
  \caption{$95\%$ Prediction Intervals for New Observations}
\end{table}

As mathematically expected, these prediction intervals are
significantly wider than their corresponding confidence
intervals. This increased width explicitly accounts for the
inherent random variation ($\sigma^2$) between individual
penguins, even when they belong to the exact same species
and share identical physical measurements.

\subsection{Hypothesis Testing}
To formally evaluate the explanatory power and structural
validity of our multiple log-linear regression model, we
conduct three distinct tiers of hypothesis testing.

\subsubsection{Overall Model Significance}
The global F-test evaluates whether the model as a whole
explains a statistically significant portion of the variance
in expected log body mass. To do this, we compare the full
model against a restricted, intercept-only baseline model.
\begin{itemize}
    \ii \textbf{Null Hypothesis ($H_0$):} $\beta_1 = \beta_2 =
    \beta_3 = \beta_4 = \beta_5 = 0$. None of the predictors
    have a significant linear relationship with the response.

    \ii \textbf{Alternative Hypothesis ($H_1$):} At least one
    $\beta_i \neq 0$ (for $i \in \{1,2,3,4,5\}$).
\end{itemize}
The test statistic
evaluates the reduction in SSE when
all predictors are added to the intercept model:
\[
  F_0 = \frac{{(\text{SSE}_{\text{restricted}} -
    \text{SSE}_{\text{full}})} / {\mathrm{df}_{\text{restricted}} -
  \mathrm{df}_{\text{full}}}}{\text{SSE}_{\text{full}}
  / \mathrm{df}_{\text{full}}}
.\]
For the global test, the restricted intercept-only model
has $\mathrm{df} = n-1$, and the full model containing $5$ predictors
has $\mathrm{df} = n-6$. Evaluating these degrees of freedom, the
formula simplifies to:
\[
  F_0 = \frac{(\text{SSE}_{\text{restricted}} -
  \text{SSE}_{\text{full}}) / 5}{\text{SSE}_{\text{full}}
  / 336}
.\]

\paragraph{Calculation and Inference}
Evaluating the global model yields an F-statistic of
$F_0 = 340.7279$, which corresponds to a p-value strictly
less than $0.001$. Because
this p-value is significantly smaller than our significance
level ($\alpha = 0.05$), we reject the null hypothesis. We
conclude with $95\%$ confidence that the regression model
as a whole explains a statistically significant proportion
of the variance in expected log body mass.

\subsubsection{Individual Parameter Significance}
While the global F-test evaluates the entire model, marginal
t-tests evaluate the independent contribution of each
specific predictor, holding all other variables constant.
\begin{itemize}
    \ii \textbf{Null Hypothesis ($H_0$):} $\beta_i = 0$. The
    $i$-th predictor does not significantly affect expected
    log body mass.

    \ii \textbf{Alternative Hypothesis ($H_1$):} $\beta_i \neq 0$.
\end{itemize}
The test statistic is calculated by dividing the parameter
estimate by its standard error:
\[
  t_0 = \frac{\hat{\beta}_i}{\text{SE}(\hat{\beta}_i)}
.\]

\begin{table}[H]
  \centering
  \begin{tabular}{|| c | c | c | c ||}
    \hline
    \textbf{Predictor Variable}
    & \textbf{Estimate ($\hat{\beta}_i$)}
    & \textbf{t-statistic}
    & \textbf{p-value} \\
    \hline\hline
    \texttt{Intercept} & 0.0759 & 0.112 & 0.911 \\
    $\log(\texttt{Flipper\_Length})$ & 0.9036 & 6.053 & $<0.001$ \\
    $\log(\texttt{Culmen\_Length})$ & 0.4364 & 5.722 & $<0.001$ \\
    $\log(\texttt{Culmen\_Depth})$ & 0.6181 & 7.682 & $<0.001$ \\
    \texttt{Is\_Gen} & 0.2341 & 6.641 & $<0.001$ \\
    \texttt{Is\_Chin} & -0.1195 & -5.925 & $<0.001$ \\
    \hline
  \end{tabular}
  \caption{Marginal t-Tests for Individual Parameters}
\end{table}

\paragraph{Inference}
As shown in the table, with the exception of the intercept,
every morphological and categorical predictor yields a
p-value strictly less than $0.001$. This strongly rejects
the null hypotheses at the $\alpha = 0.05$ level, confirming
that each variable is a highly significant independent
predictor of log body mass.

\subsubsection{Joint Significance of Categorical Variables}
To rigorously justify the inclusion of the categorical
species indicators, we test their joint significance against
a restricted model containing only the physical measurements.
\begin{itemize}
    \ii \textbf{Null Hypothesis ($H_0$):} $\beta_4 = \beta_5 = 0$.
    The species classifications provide no additional
    explanatory power beyond physical dimensions.

    \ii \textbf{Alternative Hypothesis ($H_1$):} At least one
    species parameter is non-zero.
\end{itemize}
We use the same F-statistic as the global test:
\[
  F_0 = \frac{{(\text{SSE}_{\text{restricted}} -
    \text{SSE}_{\text{full}})} / {\mathrm{df}_{\text{restricted}} -
  \mathrm{df}_{\text{full}}}}{\text{SSE}_{\text{full}}
  / \mathrm{df}_{\text{full}}}
.\]
For our specific partial test, the restricted model contains
$3$ continuous predictors ($\mathrm{df} = n-4$), and the full model
adds $2$ categorical indicators ($\mathrm{df} = n-6$). Evaluating
the degrees of freedom yields:
\[
  F_0 = \frac{(\text{SSE}_{\text{restricted}} -
  \text{SSE}_{\text{full}}) / 2}{\text{SSE}_{\text{full}}
  / 336}
.\]

\paragraph{Calculation and Inference}
Evaluating the joint addition of the species variables
yields an F-statistic of $F_0 = 95.7775$ with a
corresponding p-value strictly less than $0.001$.
Because this p-value is significantly
smaller than $\alpha = 0.05$, we firmly reject the null
hypothesis. This mathematically suggests that the categorical
species indicators provide a highly significant improvement
to the model's predictive power, reasonably justifying their
inclusion over a model based purely on physical dimensions.

\subsection{Model Diagnostics}
\label{subsection:goodness_of_fit}

To conclude our analysis, we must verify both the
explanatory power of the model and the structural validity
of the Ordinary Least Squares (OLS) assumptions in log-space.

\subsubsection{Goodness of Fit}
The explanatory power of the model is measured by the
proportion of total variance in $\log(\texttt{Body\_M})$ that is
captured by our selected predictors.
\begin{itemize}
    \ii \textbf{Multiple $R^2$ ($0.8353$):} Our model explains
    $83.53\%$ of the sample variance in the observed
    $\log(\texttt{Body\_M})$.
    \ii \textbf{Adjusted $R^2$ ($0.8328$):} Because adding
    parameters mathematically inflates the raw $R^2$. Our high
    adjusted $R^2$ suggests that the inclusion of the categorical
    species indicators provides a genuine increase in predictive
    accuracy.
\end{itemize}

\subsubsection{Residual Analysis}
A regression model is only mathematically valid if its
unobserved error terms ($\eps$) are random
and independent.

\begin{itemize}
    \ii \textbf{Residuals vs. Fitted:} The plotted residuals
    display a random scatter clustered tightly around the
    horizontal zero-line with no distinct funnels or curves.
    This suggests that the variance of the errors is constant
    and that the log-linear structural
    relationship is appears appropriate.
    % \begin{figure}[H]
    %   \centering
    %   \includegraphics[width=0.6\textwidth,keepaspectratio]{./codes/plots/model-testing/res-vs-fit.pdf}
    %   \caption{Residuals vs. Fitted}
    % \end{figure}

    \ii \textbf{Normal Q-Q Plot:} The standardized residuals
    tightly hug the theoretical $45$-degree diagonal line,
    visually indicating that the error terms are normally
    distributed.
    % \begin{figure}[H]
    %   \centering
    %   \includegraphics[width=0.6\textwidth,keepaspectratio]{./codes/plots/model-testing/qq-plot.pdf}
    %   \caption{Normal Q-Q Plot}
    % \end{figure}
\end{itemize}

\subsubsection{Normality Assessment}
To rigorously prove the visual assumption of normality
established by the Q-Q plot, we apply the Pearson Chi-Square
Goodness-of-Fit test to the model's standardized residuals.

By standardizing the residuals ($\frac{\eps - \bar{\eps}}{s}$),
we map them directly to a standard Normal distribution
curve, $N(0,1)$. We slice this theoretical curve into $k=10$
regions with equal probabilities, meaning the expected probability
of a residual falling into any specific region is exactly
$E_i = 10\%$.

We then map our actual observed standardized residuals
($O_i$) into these regions. The test statistic evaluates the
difference between the observed counts and the theoretically
expected counts:
\[
  \chi^2 = \sum_{i=1}^{k} \frac{(O_i - E_i)^2}{E_i}
.\]

% \begin{figure}[H]
%   \centering
%   \includegraphics[width=0.8\textwidth,keepaspectratio]{./codes/plots/model-testing/pearson-chisq.pdf}
%   \caption{Chi-Square Goodness-of-Fit Test for Normality}
% \end{figure}

\paragraph{Calculation and Inference}
Evaluating the standardized residuals across $10$ equiprobable
bins yields a test statistic of $\chi^2 = 9.6374$. With
$k-3 = 7$ degrees of freedom (accounting for the estimation
of the mean and variance), this corresponds to a p-value of
$0.2101$.

Because this p-value is strictly greater than our significance
threshold ($\alpha = 0.05$), we \textbf{fail to reject} the
null hypothesis. This mathematically suggests that our residuals
do not significantly deviate from a normal distribution.
Consequently, all confidence intervals, prediction intervals,
and hypothesis tests calculated in the preceding sections are
reasonably supported.

\section{Conclusion}

This study developed a multiple linear regression model to
explain penguin body mass using flipper length, culmen
length, culmen depth, and species. Motivated by allometric
scaling, a log transformation was applied to obtain a
log-linear model that better captures the underlying
relationships.

The fitted model demonstrated strong explanatory power,
accounting for about 83.5\% of the variation in log body
mass. All predictors, except the intercept, were highly
significant. Flipper length, culmen length, and culmen depth
each showed a positive association with body mass, while the
species indicators captured meaningful differences among
penguin species.

The inferential results supported the usefulness of the
model. The overall F-test confirmed statistical significance,
and the joint test showed that species provides additional
explanatory power beyond physical measurements. Confidence
and prediction intervals were consistent with the fitted
model.

Diagnostic analysis indicated that the regression assumptions
are reasonably satisfied. Residual plots showed no major
patterns, the Q-Q plot suggested approximate normality, and
the chi-square test did not detect significant deviations.
Overall, the model provides both strong predictive
performance and a meaningful representation of the
relationship between body mass and morphological traits.

\appendix
\section{Appendix}

\subsection{Dataset}

The dataset used in this study is the
Palmer Penguins dataset\cite{palmerpenguins},
originally collected by Kristen B.\ Gorman and made publicly
available by Allison Marie Horst, Alison Presmanes Hill, and
Kristen B.\ Gorman.

It can be accessed through the following sources:
\begin{itemize}
  \item GitHub:
    \href{https://github.com/allisonhorst/palmerpenguins}
    {github.com/allisonhorst/palmerpenguins}

  \item Kaggle:
    \href{https://www.kaggle.com/datasets/parulpandey/palmer-archipelago-antarctica-penguin-data}
    {kaggle.com/datasets/parulpandey/palmer-archipelago-antarctica-penguin-data}
\end{itemize}

\subsection{Code}

The code used in this project is divided based on functionality:
\begin{itemize}
  \item R Code (Computation and Model Fitting):
    \href{https://github.com/Bubu-Droid/palmer-penguins-lin-reg-sem-2/tree/main/codes/r-codes}
    {R Code Link}

  \item Python Code (Visualization and Plotting):
    \href{https://github.com/Bubu-Droid/palmer-penguins-lin-reg-sem-2/tree/main/codes/plots}
    {Python Code Link}
\end{itemize}

\bibliographystyle{plain}
\bibliography{refs}

\end{document}
